% =====================================================================
\clearpage
\part{What was proved}
% =====================================================================
\section{Proven theorems}

\figproofdag
\label{sec:theorems}
% =====================================================================

This chapter states each theorem with its precise hypothesis and label. Every theorem names the
file containing the proof. We separate the unconditional strong-coupling results from the
conditional ones and state explicitly, where relevant, that a residue theorem is \textbf{not} a
gap theorem. The proofs are given in enough detail to be checkable against the corpus; where a
standard analytic result is imported, it is named and the debt is displayed.

\subsection{The strong-coupling residue theorem (P3)}
\label{sec:p3}

\src{\code{research/unconditional/PROOF-P3-UNCONDITIONAL.md}; \code{PROOF-P1P2.md};
\code{PROOF-P3.md}; \code{P3-CYLINDER-RESIDUE.md}; \code{P3-FOURWALL-DERIVATION.md};
\code{P3-FULLREP-KP.md}; \code{P3-DRESSED-MASS.md}; \code{P3-AUDIT-CORRECTED.md}.}

\begin{claimbox}{cT}
\claimhead{Theorem}{(P3)} \emph{Fix $a>0$. There exist $\beta_1=\min(1/48,c/N)>0$ and, for
each $\beta<\beta_1$, a constant $r_0(\beta)=\tfrac12 c_1(\beta)u(\beta)^p>0$ ($c_1$ a positive
$\mathrm{SU}(N)$ group constant, $u(\beta)$ the fundamental character coefficient,
$u\asymp\beta$; the exponent $p$ is fixed by the character normalization --- the corpus uses
$p\ge1$ --- and its value only shifts the prefactor), such that, on the cylinder
$\mathcal C_L=T^3_L\times\mathbb R$ and for every spatial volume $L$,}
\begin{equation}
f_{1,L}^2:=\nu_L(\{E_{1,L}\})=\lim_{t\to\infty}e^{E_{1,L}t}\,C_L(t)\ \ge\ r_0(\beta)\ >\ 0,
\label{eq:p3}
\end{equation}
\emph{(limit $t\to\infty$ at fixed $L$ first, then $L$-uniformity; $L\to\infty$ is not taken).}
\end{claimbox}

\textbf{Hypotheses and imports.} Wilson action
$S=\beta\sum_p\bigl(1-N^{-1}\operatorname{Re}\operatorname{tr}U_p\bigr)$, Haar measure, character
expansion with $c_r(\beta)\ge0$. The proof imports the Osterwalder--Seiler connected-composite
expansion and the Koteck\'y--Preiss criterion; the ``standard-assembly'' lemmas are proved in
full in the four companion files: \code{P3-CYLINDER-RESIDUE.md} (cylinder plus spectral bridge),
\code{P3-FOURWALL-DERIVATION.md} (four-wall geometry), \code{P3-FULLREP-KP.md}
(full-representation KP, $K\le CN$), \code{P3-DRESSED-MASS.md} (dressed-mass resummation), modulo
the standard OS/K\"all\'en--Lehmann representation. The hypothesis is the per-site condition
$u(\beta)K_{\mathrm{KP}}<1$ with explicit $K_{\mathrm{KP}}\le CN$ (the plaquette adjacency
degree is $\le20$; the $\mathrm{SU}(N)$ link multiplicity is polynomial in $N$, bounded by $CN$
on the fundamental sector). The one constant not reduced to a decimal is the constant in
$C_g\le CN$; its \emph{existence} is proved, and its value only shifts the prefactor of
$\beta_1$.

\textbf{Proof skeleton.} (i) $\operatorname{tr}F^2$ is a finite A1g sum of local plaquette
observables, so $C(t)$ is a sum over connected polymers; absolute convergence is the KP condition
$uK<1$, uniform in volume and strip length. (ii) A connected polymer linking time $0$ to time $t$
must meet every intermediate slice, so its plaquette cost is $\ge 4t/a$ (a four-walled tube);
equality forces the straight four-walled tube, with weight
$a_0u^{4t/a}=a_0e^{-m_1t}$, $m_1=4\ln(1/u)/a$. This tube is the \emph{unique} leading
concentration, and its amplitude $a_0(\beta)>0$ is a product of positive fundamental character
coefficients times the positive singlet recoupling $f\otimes\bar f\to1$. (iii)
$f_1^2=|\langle\Omega|O|\mathbb{G}\rangle|^2$ is a modulus-square, so no sign can undo it.
(iv) Every non-tube polymer costs a whole factor $u$ more; the KP criterion is independent of
backbone length, so the remainder is $O(u)$ uniformly in $t,V,L$.

\textbf{What it is.} A genuine \textbf{volume- and $L$-uniform lower bound on the residue}
$f_1^2$ at strong coupling, fixed $a$. It is the half that Shen--Zhu--Zhu does not supply: SZZ
gives a gap (a decay-rate \emph{upper} bound), P3 gives a residue \emph{lower} bound. It is the
campaign's only unconditional positive statement about the atom weight.

\textbf{What it is not --- the correction that must be preserved.} Earlier drafts claimed the
integer cost grading ``additionally yields I2 (isolation).'' \textbf{That claim is wrong} and is
corrected in \code{PROOF-P3-UNCONDITIONAL.md} \S5 and \code{PROOF-P1P2.md} \S3. The cost grading
grades \textbf{polymer amplitudes}, not \textbf{transfer-matrix eigenvalue spacing}. It makes the
single tube the unique leading term (I1), but the KP/OS remainder is a $t$-uniform $O(u)$, never
an $O(e^{-\delta t})$; so it cannot separate a second atom from one at $\varepsilon\to0$, nor
place the excited $0^{++*}$ (at $\approx1.5\,m_1$) below the $2m_1$ continuum. \textbf{I2 remains
exactly the open step.} Hence P3 is a \textbf{residue theorem, not a gap theorem.} The
strong-coupling gap itself is classical Osterwalder--Seiler (1978) and is not what P3 proves.

\textbf{Scope.} Fixed $a$; $\beta<1/48$; no continuum limit $a\to0$; no transport to
$\beta\approx2.5$. The analyticity radius $1/48$ is the Ursell series radius, not a wall in the
spectrum: the residue $f_1^2(\beta)=|a(\beta)|^2$ is real-analytic wherever $\lambda_1$ is simple
and isolated (Kato; analytic in a disk), so $1/48$ is a limitation of the \emph{method}, not of
the object. But the method stops there, and no absolute strong-coupling method can cross its
analyticity radius.

\subsubsection{The polymer expansion and the Koteck\'y--Preiss criterion}

The analytic engine behind P3 is the polymer expansion, and it is worth stating in its standard
form so that the reader can see exactly what the proof imports. Let $\mathcal P$ be a countable
set of \emph{polymers} --- finite connected sets of plaquettes --- with complex activities
$w(\gamma)$, and let $\gamma\cap\gamma'\ne\emptyset$ mean that the two polymers share at least
one plaquette. The Koteck\'y--Preiss criterion is: if there is $a>0$ such that for every polymer
$\gamma_0$,
\begin{equation}
\sum_{\gamma\,:\,\gamma\cap\gamma_0\ne\emptyset}|w(\gamma)|\,e^{a|\gamma|}\ \le\ a ,
\label{eq:kp}
\end{equation}
then the sum over disjoint collections of polymers converges absolutely, the logarithm of the
partition function is a sum over connected clusters, and truncated correlations decay at rate
$a$.

In the P3 setting, the activity of a polymer is bounded by $u^{|\gamma|}$ times a combinatorial
factor, where $u(\beta)=u$ is the fundamental character coefficient ($u\asymp\beta$). The
plaquette adjacency degree is at most $20$, and the $\mathrm{SU}(N)$ link multiplicity is
polynomial in $N$, bounded by $CN$ on the fundamental sector. Hence \eqref{eq:kp} holds for
$uK<1$ with an explicit $K\le CN$, which is the per-site condition quoted in the theorem. The
constant in $C_g\le CN$ is not reduced to a decimal; its existence is what the criterion needs,
and its value only shifts the prefactor of $\beta_1$.

The four-wall geometry enters as follows. A connected polymer that links time $0$ to time $t$
must meet every intermediate time slice. Each meeting costs at least one plaquette of the tube,
and in $3+1$ dimensions a tube that links the two times has four walls, so the minimal cost is
$4t/a$ plaquettes, i.e.\ a factor $u^{4t/a}=e^{-m_1t}$ with $m_1=4\ln(1/u)/a$. The straight
tube is the unique minimizer; every other polymer costs at least one extra factor $u$. The
criterion \eqref{eq:kp} is independent of the backbone length, so the remainder is $O(u)$
uniformly in $t,V,L$. This is the entire mechanism of P3: a positive amplitude $a_0(\beta)$
times a unique exponential, with a uniform error. What it does \emph{not} give is a second
exponential scale, which is why I2 remains open.

\subsection{The no-dissolution kinematic theorem}
\label{sec:nodiss}

\src{\code{research/unconditional/PROOF-NO-DISSOLUTION-KINEMATIC.md};
\code{PROOF-DISSOLUTION-SPECTRAL-DENSITY.md}; \code{PROOF-ISOLATION.md};
\code{PROOF-EDGE-E2.md}.}

\begin{claimbox}{cC}
\claimhead{Theorem}{(no dissolution)} \emph{Suppose (K-A) the theory has a uniform gap
$m_g\ge\mu>0$, and (K-B) the lightest $0^{++}$ satisfies $m_{0^{++}}<2m_{\mathrm{lightest}}$
(implied by $0^{++}$ being the overall lightest). Then the lightest $0^{++}$ lies strictly below
the two-body ($gg$) continuum of its channel: $\theta^{2b}_{0^{++}}\ge2m_g$ and
$m_{0^{++}}<2m_g$, so the margin $2m_g-m_{0^{++}}$ is positive; it equals $m_{0^{++}}$ when the
$0^{++}$ is the overall lightest ($m_{0^{++}}=m_g$). (The operative isolation threshold is
sharper; see below.)}
\end{claimbox}

\textbf{Proof skeleton.} In the confined phase there is no asymptotic gluon, hence no $(gg)$ cut;
the lowest continuum in the $0^{++}$ channel is two glueballs, at $\theta\ge2m_g$. If the
$0^{++}$ is the overall lightest, $m_{0^{++}}=m_g$, so $m_{0^{++}}<2m_{0^{++}}\le\theta$. The atom
sits strictly below the two-body continuum by its own mass. Dissolution would require a continuum
at or below $m_{0^{++}}$, i.e.\ a state lighter than $\tfrac12m_{0^{++}}$ --- a level crossing
\textbf{from below} --- or gap closure. Same-channel $0^{++}$ versus excited $0^{++*}$ is a von
Neumann--Wigner avoided crossing (order cannot flip); a cross-channel $\{0^{++},2^{++}\}$
ordering flip is a relabeling (residues unchanged), and true dissolution needs the factor-2
hierarchy $m_{0^{++}}\ge2m_{\mathrm{lightest}}$.

\textbf{Sharpening.} In $0^{++}$, charge conjugation $C$ and parity $P$ cannot separate $gg$ from
$(gg)(gg)$ (both are $C$-even, $P$-even), so the \emph{operative} threshold $\Lambda_{\star}$ for
isolation is the excited $0^{++*}$ at $\approx1.5\,m_g$ (physically $1.56\,m_{0^{++}}$), not the
two-glueball $2m_{0^{++}}$. The coarser target ``no $0^{++}$ below $2m_g$'' is false; the correct
I2 is a gap to the excited glueball of size $\delta\approx0.5\,m_1$, fixed by power counting, not
symmetry.

\textbf{Conditional.} The theorem is a genuine kinematic theorem but inherits rather than proves
its antecedents: (K-A) is the gap, and (K-B) is the ordering input (\code{IN-11}). It fails only
on gap closure (deconfinement, not the zero-temperature smooth crossover) or on a factor-2
hierarchy. Both are dynamics, neither a smooth-crossover effect. The result is therefore \LC: it
converts ``is the $0^{++}$ an atom?'' into ``is it below the two-body threshold?'' --- and the
antecedent is exactly the open dynamics.

\subsection{Non-negativity of the Haar geometric curvature}
\label{sec:wgeo}

\src{\code{research/unconditional/PROOF-JACOBIAN-SAVES.md}.}

\begin{claimbox}{cT}
\claimhead{Theorem}{($W_{\mathrm{geo}}\ge0$)} \emph{Work in the flat chart
$X_\ell=\log U_\ell\in\mathfrak{su}(N)$, so Haar is $dU=J\,dX$ with $\Phi=-\log J$. For
$\mathrm{SU}(N)$, $\Phi=\sum_{i<j}g(\lambda_i-\lambda_j)$ with
$g(d)=-2\log\bigl(\sin(d/2)/(d/2)\bigr)$; on $|d|<2\pi$, $g''\ge1/6$. Hence $\Phi$ is convex,
$e^{-\Phi}dX$ is log-concave, and by Pr\'ekopa every linear marginal $\xi=\langle v,\log U\rangle$
has convex free energy:}
\begin{equation}
W_{\mathrm{geo}}:=E[\Phi'']-\operatorname{Var}(\Phi')\ \ge\ 0
\qquad\text{for every mode, every } k \text{ including } k\to0,\text{ uniformly in } L.
\label{eq:wgeo}
\end{equation}
\emph{Additionally, pointwise, $\Phi''\succeq0$ under any measure and any $\beta$, giving
Lemma B: $E[\Phi''|0]\ge(N/12)\|v\|^2$ ($\mathrm{SU}(2)$: $1/6$; $\mathrm{SU}(3)$: $1/4$;
measured $0.34$, $0.46$). At $\beta=0$ the geometric term is the whole curvature: $\mathrm{SU}(2)$
$V_{\mathrm{eff}}''=0.248$ versus $E[\Phi'']-\operatorname{Var}(\Phi')=0.246$, an exact
closure.}
\end{claimbox}

\textbf{The exact split and what fails.} The identity is
\begin{equation}
\begin{aligned}
G''(0)&=E\bigl[H_{vv}+\Phi''\,\big|\,\xi=0\bigr]
        -\operatorname{Var}\bigl(G_v+\Phi'\,\big|\,\xi=0\bigr)
        =\Delta_{\mathrm{dyn}}+W_{\mathrm{geo}}+X,\\
\Delta_{\mathrm{dyn}}&=E[H]-\operatorname{Var}(G),\qquad
X=-2\operatorname{Cov}(G,\Phi').
\end{aligned}
\label{eq:split}
\end{equation}
$W_{\mathrm{geo}}$ is the curvature of the pure-Haar marginal; $X$ is the cross-fluctuation of
the action force with the Jacobian. The theorem bounds $W_{\mathrm{geo}}$ from \textbf{below by
0}, but at $\beta=2.5$ the measured decomposition attributes the sign to this term
($W_{\mathrm{geo}}=+0.319\pm0.005$ on the probe of Chapter~\ref{ch:numerics}, \LE; that probe is
gauge-dependent, Appendix~\ref{app:postaudit}). Certifying
$W_{\mathrm{geo}}>\operatorname{Var}-E[H]$ needs an upper bound on $\operatorname{Var}(G|0)$,
which is precisely the induced-mass/H1 object. So this is a \textbf{sign theorem, value-free},
not a gap theorem. The $-0.165$ ``DEAD-lean'' value reported elsewhere is a proxy that dropped
exactly $+W_{\mathrm{geo}}$; measured exactly on its own mode, the marginal is $+0.51$.

\subsection{Gauge-invariant Perron--Frobenius, and the edge no-go}

\src{\code{research/unconditional/PROOF-GAUSS-INVARIANT-PF.md}.}

\begin{claimbox}{cT}
\claimhead{Theorem}{(gauge-invariant PF)} \emph{On a finite lattice, $\beta>0$, the transfer
matrix $T=e^{-aH}$ has a strictly positive kernel $T(U,V)>0$ for all $U,V$ (a limit of the
electric heat kernel, which is positive by
$K_t(g,h)=\sum_\lambda d_\lambda\chi_\lambda(g^{-1}h)e^{-tC_\lambda}>0$, times a positive
magnetic weight). Since $T(gU,gV)=T(U,V)$, it descends to a strictly positive kernel
$\bar T$ on $\mathcal{A}/\mathcal{G}$; for $0\le f\in \mathcal{H}_{\mathrm{phys}}$, $f\ne0$,
$(Tf)(U)>0$ everywhere. So $T$ is positivity-improving at every $\beta>0$, and by
Jentzsch/Krein--Rutman $\lambda_0=1$ is simple with a strictly positive vacuum $\Omega\equiv1$,
and $0\le\lambda_1<\lambda_0$, giving a finite-volume gap
$\Delta_V=-(1/a)\log\lambda_1>0$.}
\end{claimbox}
\label{thm:pf}

\textbf{Why this is honest and why it falls short.} The proof uses the \textbf{full positive Haar
measure}, no gauge slice, so no Gribov copy or Faddeev--Popov sign enters --- the honest content
of ``a positivity gauge-fixed PF lacked.'' But (i) $\lambda_1$ is the lightest state over
\emph{all} channels, not the $0^{++}$ specifically; (ii) $\inf_V\Delta_V>0$ is the mass gap and
is not obtained --- PF/Dobrushin constants scale like $e^{-cV}$, so finite-volume gaps are
\emph{upper} bounds on $\Delta_{\mathbb{R}^4}$; (iii) at $a\to0$ compactness fails, and the free
field (reflection-positive, PF, strictly positive vacuum) has $\Delta_V\to0$.

\textbf{The exact edge obstruction \LN.} The vacuum is $0^{++}$ and the lightest scalar
glueball is \emph{also} $0^{++}$: they share the channel. The $0^{++}$-projected transfer matrix
is still positivity-improving, so its Perron vector is the \textbf{vacuum} (top of the sector)
and the glueball is the second state, in $\Omega^\perp\cap(0^{++})$ --- exactly where the single
reflection-positivity cone degenerates. To reach the glueball one would need a second invariant
cone / total positivity (a Karlin flag) on $\mathcal{A}/\mathcal{G}$, but the quotient carries
\textbf{no order}, so no Gantmacher--Krein/TP cone exists. Gauge invariance reaches the vacuum
and the gap, never the glueball. The decisive check is on the physical space itself: free
Maxwell is PF and strictly positive with a finite-volume gap, yet its $0^{++}$ spectrum is a
two-photon cut from zero, $f_1^2=\nu(\{0\})=0$. Gauge-invariant PF plus strict positivity plus a
finite-volume gap \textbf{does not imply} a glueball atom.

\subsection{The strong-coupling coarse Poincar\'e lemma (Lemma SC)}

\src{\code{research/unconditional/PROOF-COARSE-POINCARE.md} \S(c).}

\begin{claimbox}{cT}
\claimhead{Lemma}{(SC)} \emph{At $\beta<1/48$, Shen--Zhu--Zhu give a uniform LSI
$\operatorname{Var}_\mu(F)\le(1/\lambda_\mu)\int|\nabla F|^2d\mu$, uniformly in volume. For a
block averaging map $P$ with $\operatorname{Lip}(P)=C_{\mathrm{blk}}=O(1)$ (averaging is a
contraction, volume-independent), set $F=\varphi\circ P$. Then for the pushed-forward measure
$\nu=P_*\mu$,}
\begin{equation}
\operatorname{Var}_\nu(\varphi)\ \le\ \frac{C_{\mathrm{blk}}^2}{\lambda_\mu}
\int|\nabla\varphi|^2\,d\nu ,
\label{eq:sc}
\end{equation}
\emph{uniformly in $(L,a)$. So the coarse Poincar\'e inequality holds at strong coupling with
$\lambda=C_{\mathrm{blk}}^{-2}\lambda_\mu$.}
\end{claimbox}

\textbf{Scope.} This is a genuine anchor, but it lives exactly where the full LSI already holds.
It does not transport to $\beta=2.5$, which lies outside every high-temperature regime. The
file's three direct routes to $\beta=2.5$ fail for independent, nameable reasons: the effective
action is Savvidy-indefinite (not provably convex; the bare coarse Hessian minimum is
$\approx-0.3$, and the one-loop Schur form is undefined because $H_{ff}$ is indefinite); the
block gap is a \emph{fast} relaxation rate, not the slow stiffness, and the only variance bound
runs the wrong way by $\approx e^{5N_p}$; and the provable regime is $\beta<1/48$.

\subsection{The Feshbach--Bogoliubov lemma and the exact variance decomposition}
\label{sec:feshbach}

\src{\code{research/unconditional/FORGE-BLIND.md} \S2; \code{INVENT-VARIANCE-INEQ.md};
\code{D-BBS-INTEGRATE-NEG.md}; \code{docs/new-tools-registry.md} A2.}

\begin{claimbox}{cT}
\claimhead{Lemma}{(two-block gap; unconditional, abstract)} \emph{Let}
\begin{equation}
\mathbb{A}=\begin{pmatrix} A & B\\ B^* & C\end{pmatrix}
\quad\text{on}\quad \mathcal{H}=\mathcal{H}_-\oplus\mathcal{H}_+,
\qquad C\ge\gamma I .
\label{eq:twoblock}
\end{equation}
\emph{Define the Schur complement $\Sigma:=A-BC^{-1}B^*$ on $\mathcal{H}_-$. If $\Sigma\ge\mu I$,
then}
\begin{equation}
\mathbb{A}\ \ge\ \min\Bigl\{\frac{\mu}{1+2\|B\|^2/\gamma^2},\ \frac{\gamma}{2}\Bigr\}\cdot I
\ >\ 0 .
\label{eq:twoblock-bound}
\end{equation}
\end{claimbox}
\begin{proofx}
$\langle u,\mathbb{A}u\rangle=\langle u_1,Su_1\rangle+\gamma\|u_2+\gamma^{-1}B^*u_1\|^2$; put
$v_2=u_2+\gamma^{-1}B^*u_1$ and use
$\|u_2\|^2\le2\|v_2\|^2+2\|B\|^2\gamma^{-2}\|u_1\|^2$.
\end{proofx}

The lemma needs \textbf{no hypothesis of a positive Hessian}; $A$ may be arbitrarily negative.
The indefinite Savvidy Hessian is not hidden but \emph{converted} into the demand $\Sigma\ge\mu>0$.
For 4D pure $\mathrm{SU}(N)$, this becomes the \textbf{tachyonic-compensation inequality} (TCI):
\begin{equation}
\langle\varphi,Q\varphi\rangle+\langle\varphi,\mathbb{K}\varphi\rangle-cF\ \ge\ \mu>0
\qquad\text{for all }\|\varphi\|=1,\ \varphi\in\mathcal{H}_- ,
\label{eq:tci}
\end{equation}
where $\mathbb{K}$ is the lift of the block scheme; it is unrelated to the fibre Hessian $K$ of
the variance decomposition below. Inequality \eqref{eq:tci} says that the gauge-invariant
effective potential is convex at its minimum, pointwise and nonperturbatively, and the Savvidy
vacuum has no imaginary part. The tool is not vacuous: on the free massless field, $Q=0$, the
bottom of $A|_-$ is $0$, $\Sigma$ is not $\ge\mu$, TCI fails, and the certificate correctly
returns $0$.

\textbf{Lemma NS (eliminating a negative block raises the survivor).} In the quadratic form
$\begin{psmallmatrix}A&B\\B^*&C\end{psmallmatrix}$ with $C\prec0$, the Schur complement satisfies
$\Sigma=A-BC^{-1}B^*=A+B|C|^{-1}B^*\succeq A$. This is true and is a genuine gain, but it
certifies only the \textbf{fast} half of the block BBS scheme; the slow wall is the variance term,
which the block elimination never touches.

\textbf{Exact variance decomposition (Savvidy-corrected Brascamp--Lieb).} Here $V$ is the full
action and $H:=\operatorname{Hess}V$ its Hessian --- not the Hamiltonian $H_a$ of
Chapter~\ref{sec:intro}. With $K=\Pi_0\operatorname{Hess}V\,\Pi_0$,
$\mathbf{c}=\Pi_0\operatorname{Hess}V\,v$, $h_n=v^{\mathsf T}\operatorname{Hess}V\,v$,
\begin{equation}
\begin{aligned}
G''(0)&=S(v)-D(v),\\
S&=E_\mu\bigl[h_n-\langle K^{-1}\mathbf{c},\mathbf{c}\rangle\mid0\bigr],\qquad
D=\operatorname{Var}_\mu(\partial_nV\mid0)-E_\mu\bigl[\langle K^{-1}\mathbf{c},\mathbf{c}\rangle\mid0\bigr].
\end{aligned}
\label{eq:vardecomp}
\end{equation}
Sign law: $D\le0\iff K\succeq0$ (true BL); $D=0\iff$ Gaussian; $D>0\iff$ fiber Savvidy. At
$\beta=0$, $V=\Phi$ is convex, so $D\le0$ and $G''\ge S\ge N/12>0$ --- a theorem. \LE{} At
$\beta=2.5$ a measured $D/S\approx2\%$ was reported, but \emph{on the gauge-dependent gluon
mode}, not the gauge-invariant $0^{++}$ (Appendix~\ref{app:postaudit}); its one-sided interval
still includes $0$, and a correct-object audit finds the composite variance $V$-suppressed, so
this is a Haar/$W_{\mathrm{geo}}$ number, not a positive-gap proof. The residual --- a uniform
bound on the fiber defect $D$ --- is H1. The acid test is non-trivial: the free Gaussian has
$D=0$, and the $k\to0$ mode has $\operatorname{Var}(\xi)=\infty$, so $S=0$ and
$G''=0-0=0$; the tool manufactures no false positivity.

\subsection{The residue--level-gap lemma and the conditional theorem}
\label{sec:reslevel}

\src{\code{GAP-G2-edge} (quoted in \code{G2-ASSEMBLE-100.md} \S2);
\code{docs/conditional-theorem.md}.}

\begin{claimbox}{cT}
\claimhead{Lemma}{(residue--level-gap)} \emph{Suppose the lightest state is an isolated atom at
$m_1^2$, and let $\Lambda_{\star}:=\inf\operatorname{supp}(\nu-f_1^2\delta_{m_1^2})>m_1^2$.
Then}
\begin{equation}
f_1^2\ \ge\ \frac{\Lambda_{\star} M_0-M_1}{\Lambda_{\star}-m_1^2}
= M_0\,\frac{\Lambda_{\star}-\bar R}{\Lambda_{\star}-m_1^2},
\qquad \bar R:=\frac{M_1}{M_0}.
\label{eq:reslevel}
\end{equation}
\emph{The right side is positive \textbf{iff} $\Lambda_{\star}>\bar R$, and $\Lambda_{\star}>\bar
R$ suffices (if $f_1^2=0$ then $\Lambda_{\star}=m_1^2$ and $\bar R\ge\Lambda_{\star}$). The
converse fails (a positive atom with a heavy tail can have $\bar R>\Lambda_{\star}$). So this is
a self-certifying sufficient certificate, not an equivalence. By this lemma and Theorem A of
\code{mass-spectrum}, $\text{atom(I1)}\wedge\text{I2}\iff$ sharp two-exponential decay
$C(t)=f_1^2e^{-m_1t}(1+\eta(t))$, $|\eta(t)|\le Be^{-\delta t}$.}
\end{claimbox}

\begin{proofx}
Write $\nu=f_1^2\delta_{m_1^2}+\mu$, $\mu\ge0$, $\operatorname{supp}\mu\subset
[\Lambda_{\star},\infty)$. Then $M_1=f_1^2m_1^2+\int s\,d\mu$ and
$M_0=f_1^2+\int d\mu$, so $\int d\mu\le M_1/\Lambda_{\star}$ and hence
$M_0\ge f_1^2+M_1/\Lambda_{\star}$, i.e.\ $f_1^2\le M_0-M_1/\Lambda_{\star}$ (an upper bound,
not the claim). For the lower bound use the two-point extremal: among measures with the given
$M_0,M_1$ supported in $\{m_1^2\}\cup[\Lambda_{\star},\infty)$, the minimal atom weight is
attained when all the continuum mass sits at $\Lambda_{\star}$, giving
$f_1^2=(\Lambda_{\star}M_0-M_1)/(\Lambda_{\star}-m_1^2)$. If the actual continuum sits higher,
the atom is larger. This proves the displayed bound.
\end{proofx}

\textbf{The conditional theorem.} Under the twelve named inputs \code{IN-0} \ldots{} \code{IN-11}
(\code{docs/conditional-theorem.md}), there exists a gauge-invariant measure $\mu$ on
$\mathbb{R}^4$ satisfying OS0--OS4, non-Gaussian with $g_{\mathrm{ren}}/\Lambda>0$, with
$\Delta/\Lambda\ge\min(\delta_W,\delta_S)>0$, unique vacuum, area law, and group-uniform
constants; and the gap is equivalent to
$\inf\operatorname{spec}(H_a|_{\Omega^\perp})\ge\delta\Lambda$ uniformly in $a$. The synthesis
chain is \textbf{9/9 closed} (each gluing lemma is either proved or reduced to a named input),
and the core of S3, S4, S5, S10, S11 is machine-verified in Lean with frozen axioms
(\code{Handoff.lean}, \code{RGHandoff.lean}, \code{HandoffMargin.lean},
\code{GlueA.spectral\_decay}, \code{GL2a.lean}, \code{GL2bPrime.lean}, \code{GL7Dual.lean},
\code{GroupUniform.lean}, \code{GL6Fubini.lean}, \code{Conditional.lean}; the checks are listed
in \code{conditional-theorem.md} \S6).

\textbf{Honest status of the inputs.} Twelve inputs: two proved or machine-verified
(\code{IN-1} bounded class, vacuous, and \code{IN-8} assembly), seven open (\code{IN-2},
\code{IN-3}, the 4D part of \code{IN-4}, \code{IN-5}, \code{IN-6}, the quantitative constants of
\code{IN-7}, \code{IN-9}), one explicit magnitude bridge (\code{IN-10}, \texttt{hcert}:
$\text{cert}\le\Delta/\Lambda$), and one physical input (\code{IN-11}, the lightest bulk state in
$0^{++}/2^{++}$). \code{IN-3} --- the nonperturbative beta-function bound in the
intermediate-coupling region $\lambda_{\mathrm{tH}}\in[1,7.3]$ --- is \textbf{the single physical
wall}. Balaban-type RG stops at $\lambda\sim1$; strong-coupling expansions require $\lambda$
large; there is no rigorous method in between. The theorem is \LC, and the conditional reads as:
\emph{if} these twelve inputs hold, Clay follows --- not that they do.

\textbf{A crucial correction (red team).} An earlier campaign summary overstated ``the
certificate is the gap.'' The correct statement (\code{campaign-final.md} \S6) is that the
certificate is a \textbf{sufficient} route
($\text{cert}+\text{IN-10}+\text{IN-11}\Rightarrow\text{gap}$), but $\text{cert}>0$ is strictly
weaker than the gap, and the necessity $\text{gap}\Rightarrow\text{cert}$ is unproved. So
``stuck at the certificate'' does \textbf{not} imply ``stuck at the gap''; routes to the gap
outside the certificate family are not ruled out. This paper preserves that correction.

\subsection{The strong-coupling bridge and the conditional roof}
\label{sec:roof}

Two further results arose in the roof wave and are worth recording as theorems with explicit
scope.

\begin{claimbox}{cC}
\claimhead{Theorem}{(strong-coupling bridge)}
\emph{On the strong-coupling side, the area law and the $0^{++}$ gap share the same polymer
weight $u$; the leading tube gives $m_1=4\ln(1/u)/a$ and the string tension $\sigma$
gives $m_1/\sqrt\sigma=4\sqrt{\ln(1/u)}\ge4$. Hence}
\[
m_{0^{++}}\ \ge\ 4\sqrt{\sigma}
\]
\emph{for $\beta<\beta_*$ with $\beta_*\le8.9\times10^{-3}$.}
\end{claimbox}

The constant is explicit and the inequality is genuine, but the regime is far inside strong
coupling. It is the campaign's one place where confinement is \emph{proved} to imply a gap with a
numerical constant.

\begin{claimbox}{cT}
\claimhead{Theorem}{(conditional roof)}
\emph{Let $\chi$ be the certified susceptibility, $\gamma$ the Cheeger constant of the coarse
probability measure, and $h$ its Cheeger profile. In the family where the profile is
$h\sqrt\chi=\sqrt{2/\pi}$ (exact at every $\sigma$ in the Gaussian coarse model), one has}
\[
\gamma\ \ge\ \frac{1}{2\pi}\,\frac{1}{\chi_{\mathrm{certified}}} .
\]
\end{claimbox}

Given the $\chi$-certificate, the roof is laid at a rigorous fraction of about $4.5\%$. The
$\chi$-certificate is the one missing strut; the theorem is the precise sense in which the roof
is \emph{conditionally} layable. Figure~\ref{fig:roof} draws the three names of the strut.

\figroof

\subsection{Summary table of proved theorems}

\begin{table}[t]
\centering
\footnotesize
\begin{tabularx}{\textwidth}{@{}>{\raggedright\arraybackslash}p{2.6cm}cc
  >{\raggedright\arraybackslash}X>{\raggedright\arraybackslash}X@{}}
\toprule
\textbf{Theorem} & \textbf{Label} & \textbf{Hypothesis} & \textbf{What it gives} &
\textbf{What it does \emph{not} give}\\
\midrule
P3 residue (I1) & \LT{} ($+$OS/KP) & fixed $a$, $\beta<1/48$, $T^4_L$ &
  $f_1^2\ge r_0(\beta)>0$, uniform in $V,L$ & I2 isolation; $a\to0$; $\beta\approx2.5$\\
No dissolution & \LC & uniform gap $+$ $0^{++}$ lightest &
  $m_{0^{++}}<2m_g$ (two-body); atom pinned & the gap $+$ ordering themselves\\
$W_{\mathrm{geo}}\ge0$ & \LT & Haar, flat chart, $g''\ge1/6$ &
  sign of the geometric term, all $k$, all $L$ & dominance over Savvidy at 2.5\\
$E[\Phi'']\ge N/12$ & \LT & pointwise, any measure & floor on the Jacobian curvature &
  strict gap\\
Gauge-invariant PF & \LT{} finite $V$ & $\beta>0$, positive Haar pushforward &
  vacuum, finite-volume gap & first excited $0^{++}$ (edge)\\
Lemma SC & \LT & $\beta<1/48$, SZZ LSI, 1-Lipschitz blocks & coarse Poincar\'e,
  $(L,a)$-uniform & transport to 2.5\\
Feshbach--Bogoliubov & \LT{} abstract & two-block, $C\ge\gamma I$, $\Sigma\ge\mu I$ &
  explicit gap, no positive Hessian needed & TCI (the 4D input)\\
$G''=S-D$ & \LT{} identity & exact & $\beta=0$ positivity $G''\ge N/12$ &
  uniform bound on $D$ at 2.5 ($=$ H1)\\
Residue--level-gap & \LT & moments available & self-certifying sufficiency
  $\Lambda_{\star}>\bar R$ & equivalence; $\Lambda$ lower bound\\
Strong-coupling bridge & \LC & $\beta<8.9\times10^{-3}$ & $m\ge4\sqrt\sigma$ &
  the crossover\\
Conditional roof & \LT{}/given & certified $\chi$, Gaussian profile &
  $\gamma\ge(1/2\pi)(1/\chi)$ & the $\chi$-certificate\\
Conditional theorem & \LC & \code{IN-0} \ldots{} \code{IN-11} & Clay conclusion &
  the inputs (esp.\ \code{IN-3})\\
Lean H1 sharpening & \LT{} (logic) & order-algebraic hypotheses & H1 $\not\Rightarrow$ atom;
  $m\ge\lambda f$ & the YM input $f>0$\\
\bottomrule
\end{tabularx}
\caption{Proved theorems and their exact scope.}
\label{tab:theorems}
\end{table}
